\section*{Chapter \ref{ch:g9:grids-and-images} --- Grids and Images}

\begin{solution}{exo:g9:grids-and-images:1}
\code{grid[2][0]} is 5, the first \omterm{def:g6:variables-and-input:value}{value} of the last row; \code{grid[0][3]} is
6, the last \omterm{def:g6:variables-and-input:value}{value} of the first row.
\end{solution}

\begin{solution}{exo:g9:grids-and-images:2}
\code{[9, 4, 4]}: the cell is the bottom right corner, so only three cells touch
it, \code{grid[1][2]}, \code{grid[1][3]} and \code{grid[2][2]}.
\end{solution}

\begin{solution}{exo:g9:grids-and-images:3}
Rows \code{255 155} and \code{55 0}: black and white swap, and each grey goes to
its opposite.
\end{solution}

\begin{solution}{exo:g9:grids-and-images:4}
Column \code{c} of the old \omterm{def:g9:grids-and-images:grid}{grid} becomes row \code{c} of the new one:

\omcode{code/grade-9/03-grids-and-images/exercises/transpose.py}{1}{9}{Rows and
columns swapped.}\label{lst:g9:grids-and-images:sol-transpose}
\end{solution}

\begin{solution}{exo:g9:grids-and-images:5}
For the bottom right cell, \code{grid[r + 1]} asks for a row that does not
exist. For the top left one, \code{grid[r - 1]} is \code{grid[-1]}: a negative
\omterm{def:g7:lists:list}{index} counts from the end, so Python quietly reads the last row, and
\code{grid[0][-1]} the last column. The fix keeps only the places inside the
\omterm{def:g9:grids-and-images:grid}{grid}:

\omcode{code/grade-9/03-grids-and-images/exercises/neighbours_fixed.py}{1}{9}{Four
\omterm{def:g9:grids-and-images:neighbour}{neighbours}, inside the \omterm{def:g9:grids-and-images:grid}{grid} only.}\label{lst:g9:grids-and-images:sol-neighbours}
\end{solution}

\begin{solution}{exo:g9:grids-and-images:6}
\code{image[::-1]} reverses the order of the rows: the image is turned upside
down. The arrow still points right, one row higher, and the grey square moves
to the bottom. A mirror reverses each row, as in
\cref{lst:g9:grids-and-images:mirror}.
\end{solution}

\begin{solution}{exo:g9:grids-and-images:7}
\omcode{code/grade-9/03-grids-and-images/exercises/brighter.py}{1}{14}{Brighter,
kept between 0 and 255.}\label{lst:g9:grids-and-images:sol-brighter}

A negative \code{amount} makes the image darker.
\end{solution}

\begin{solution}{exo:g9:grids-and-images:8}
The first row of the result is the first column, read from the bottom up:

\omcode{code/grade-9/03-grids-and-images/exercises/rotate.py}{1}{10}{A quarter
turn clockwise.}\label{lst:g9:grids-and-images:sol-rotate}

Width and height swap: \code{shape.pgm} becomes 12 \omterm{def:g5:text-pictures-sound-as-numbers:pixel}{pixels} wide and 16 high.
\end{solution}

\begin{solution}{exo:g9:grids-and-images:9}
Count the pairs of cells side by side: in each of the 4 rows, 4 pairs; in each
of the 5 columns, 3 pairs; 31 pairs in all. Each pair counts once for each of
its two cells: $2 \times 31 = 62$.
\end{solution}

\begin{solution}{exo:g9:grids-and-images:10}
\omcode{code/grade-9/03-grids-and-images/exercises/count_white.py}{1}{6}{White
\omterm{def:g5:text-pictures-sound-as-numbers:pixel}{pixels}.}\label{lst:g9:grids-and-images:sol-white}

The threshold keeps the 46 \omterm{def:g5:text-pictures-sound-as-numbers:pixel}{pixels} of the arrow, at 230; the grey square, at
120, becomes black.
\end{solution}

\begin{solution}{exo:g9:grids-and-images:11}
At most 4 \omterm{def:g7:strings:string}{characters} per \omterm{def:g5:text-pictures-sound-as-numbers:pixel}{pixel}, so about $4wh$ \omterm{def:g4:memory-and-storage:byte}{bytes}, plus a dozen for the
header. One \omterm{def:g4:memory-and-storage:byte}{byte} per \omterm{def:g5:text-pictures-sound-as-numbers:pixel}{pixel} would take $wh$: the text format is about four
times larger, the price of being readable.
\end{solution}

\begin{solution}{exo:g9:grids-and-images:12}
The new \omterm{def:g9:grids-and-images:grid}{grid} is built apart from the old one, as for the blur:

\omcode{code/grade-9/03-grids-and-images/exercises/life.py}{1}{23}{One
generation.}\label{lst:g9:grids-and-images:sol-life}

Three 1s in a row become three 1s in a column, then a row again: the middle
cell keeps 2 \omterm{def:g9:grids-and-images:neighbour}{neighbours}, the two ends have only 1 and die, and the cells above
and below the middle have exactly 3 and come to life. A test of the chapter
checks both steps.
\end{solution}

\begin{solution}{pb:g9:grids-and-images:1}
\textbf{1.} The 9 \omterm{def:g5:text-pictures-sound-as-numbers:pixel}{pixels} of the 3 by 3 square around the middle, each at
$255 / 9 \approx 28.33$.

\textbf{2.} Two \omterm{def:g5:text-pictures-sound-as-numbers:pixel}{pixels} from the middle in every direction: a 5 by 5 square.

\textbf{3.} Each blur adds a ring one \omterm{def:g5:text-pictures-sound-as-numbers:pixel}{pixel} wide: $(2k + 1) \times (2k + 1)$.

\textbf{4.} The \omterm{def:g5:text-pictures-sound-as-numbers:pixel}{pixels} that are not 0; \code{whole} rounds each new \omterm{def:g6:variables-and-input:value}{value} to a
whole \omterm{def:g5:text-pictures-sound-as-numbers:grey}{grey level}, as a \omterm{def:g9:grids-and-images:pgm}{PGM file} must.

\textbf{5.} Yes: 9, 25, 49, 81, 121, 169, the squares of 3, 5, 7, 9, 11, 13.

\omoutput{code/grade-9/03-grids-and-images/outputs/spread.txt}{Decimal \omterm{def:g6:variables-and-input:value}{values},
then whole \omterm{def:g5:text-pictures-sound-as-numbers:grey}{grey levels}.}\label{out:g9:grids-and-images:spread}

\textbf{6.} The image is only 15 by 15: after 7 blurs, the light reaches the
edges, and there are no more \omterm{def:g5:text-pictures-sound-as-numbers:pixel}{pixels} to light.

\textbf{7.} Far from the edges, yes: each \omterm{def:g5:text-pictures-sound-as-numbers:pixel}{pixel} hands out one ninth of its \omterm{def:g6:variables-and-input:value}{value}
to each of 9 \omterm{def:g5:text-pictures-sound-as-numbers:pixel}{pixels}, so the total stays 255. Once the light reaches the edges,
the share that would go outside is lost.

\textbf{8.} The same 9 and 25, then fewer: 45, 61, 69, and 69 again.

\textbf{9.} At the border of the lit area, a dark \omterm{def:g5:text-pictures-sound-as-numbers:pixel}{pixel}'s window holds a few
small \omterm{def:g6:variables-and-input:value}{values}. When their sum is less than 4.5, the average is less than 0.5,
and it rounds to 0: the \omterm{def:g5:text-pictures-sound-as-numbers:pixel}{pixel} stays dark.

\textbf{10.} When every \omterm{def:g5:text-pictures-sound-as-numbers:pixel}{pixel}'s window averages to a \omterm{def:g6:variables-and-input:value}{value} that rounds back to
the \omterm{def:g5:text-pictures-sound-as-numbers:pixel}{pixel}'s own \omterm{def:g6:variables-and-input:value}{value}, nothing changes, and every later blur gives the same
image.

\textbf{11.} Rounding errors can stop a process, or keep it where it is,
although the exact calculation would go on changing.

\textbf{12.} \textbf{Result.} With decimal \omterm{def:g6:variables-and-input:value}{values}, after $k$ blurs exactly
$(2k + 1) \times (2k + 1)$ \omterm{def:g5:text-pictures-sound-as-numbers:pixel}{pixels} are lit, away from the edges. With whole \omterm{def:g5:text-pictures-sound-as-numbers:grey}{grey
levels}, the light stops spreading at 69 \omterm{def:g5:text-pictures-sound-as-numbers:pixel}{pixels} after 5 blurs, and from blur 16
on the image no longer changes at all, its centre stuck at 3.
\end{solution}
